%=======================================================================

% Template for Thesis at Computer Sciencte at
% Ostfalia University of Applied Science
% Edit: Dirk J. Lehmann, Kai Michael Blum

%=======================================================================
\documentclass[
		11pt,
		a4paper, 
		oneside, % twoside use two sided for book layout
		% openany,
		toc=listofnumbered,
		toc=bibliography, 
		toftofoc,
		headings=small,
		headings=twolineappendix
]{scrbook}
\usepackage[utf8]{inputenc}

%=======================================================================

\newif\ifdraft
\draftfalse % Sagt aus, dass dieses Dokument ein Entwurf ist. Somit wird todonotes aktiviert. Zum deaktivieren diese Zeile auskommentieren oder auf \draftfalse setzen.

%=======================================================================

\input{Meta/praeampel}

%=======================================================================

\newcommand{\documenttitle}{TITLE OF THE THESIS}
\newcommand{\documentname}{NAME OF AUTHOR}
\newcommand{\documentmanr}{MATRIKELNUMMER}
\newcommand{\documentmodul}{}
\newcommand{\documentstudiengang}{STUDIENGANG}
\newcommand{\documentachievement}{zur Erlangung des akademischen Grades:\\Master of Science\\}




\newcommand{\documentbetreuer}{Prof. Dr.-Ing. habil. Dirk Joachim Lehmann}
\newcommand{\documentbetreuerzwei}{HIER ZWEITGUTACHTER (FALLS VORHANDEN)}
\newcommand{\documentsemester}{2022}

\newcounter{documentfakultaet}
\setcounter{documentfakultaet}{1} % 0 = Bau Wasser Boden, 1 = Informatik

%=======================================================================


\begin{document}	

    % = = = = =	= = = = = = = = = = = = = = = = = = = = = = = =
	
	\input{Kapitel/0_0_Deckblatt}
    
    % = = = = =	= = = = = = = = = = = = = = = = = = = = = = = =	
	
	\frontmatter
	\pagenumbering{Roman}
	\ifnum\value{documentfakultaet}=1
		\input{Kapitel/0_erklaerung}
		\markboth{Erklärung}{}
	\fi
    
    % = = = = =	= = = = = = = = = = = = = = = = = = = = = = = =
	
	\input{Kapitel/0_Abstract}
    
    % = = = = =	= = = = = = = = = = = = = = = = = = = = = = = =	
	
	% disable header
	\fancyhead[L]{\chaptermark}  % L -> Left part
	\fancyhead[RO]{\chaptermark} % RO -> Right part on Odd pages
	
	\tableofcontents % Inhaltsverzeichnis
	
	% = = = = =	= = = = = = = = = = = = = = = = = = = = = = = =
	
	\input{Kapitel/0_Abkuerzungsverzeichnis}
	
	% = = = = =	= = = = = = = = = = = = = = = = = = = = = = = =
	
	\mainmatter
	
	% reanable header
	\fancyhf{}
	\rhead{\nouppercase\rightmark}
	\lhead{\nouppercase\leftmark}
	
	\fancypagestyle{plain}{}
	%\fancyhead[LE]{\chaptermark} % Chapter in header Left
	%\fancyhead[RE]{\chaptermark} % Page number in header Right
	%\fancyfoot[LE]{\thepage}     % LE -> Left part on Even pages
	\fancyfoot[RO]{\thepage}     % RO -> Right part on Odd pages
	
	% = = = = =	= = = = = = = = = = = = = = = = = = = = = = = =
	
	% reset list of used acronyms
	\acresetall
	
	% = = = = =	= = = = = = = = = = = = = = = = = = = = = = = =
	
	% Kapitel einfügen
	\input{Kapitel/1_Introduction}
	\input{Kapitel/2_Related_work}
	\input{Kapitel/3_Background}
	\input{Kapitel/4_Approach}
	\input{Kapitel/5_Implementation}
	\input{Kapitel/6_Evaluation}
	\input{Kapitel/7_Discussion}
	\input{Kapitel/8_Future_work}

    % = = = = =	= = = = = = = = = = = = = = = = = = = = = = = =

    % Bilbliography
	\bibliographystyle{dinat} 
	\bibliography{literatur}
	
	% = = = = =	= = = = = = = = = = = = = = = = = = = = = = = =
	
	\ifnum\value{documentfakultaet}=0
		\input{erklaerung}
		\markboth{Erklärung}{}
	\fi
	
	% = = = = =	= = = = = = = = = = = = = = = = = = = = = = = =
	
    %\addtocontents{toc}{\cftpagenumbersoff{chapter}} % Keine Seitenzahl im Inhaltsverzeichnis
	\begin{appendices}              % Anhang Titel-Seite
	\end{appendices}
	
	\fancyfoot[R] {\thepage}        % Seitenzahl Right
	
	\appendix
	\input{Kapitel/100_Anhang}
	%\postappendix
	
	% = = = = =	= = = = = = = = = = = = = = = = = = = = = = = =
	
\end{document}